\clearpage\SourceLeaf{365}\ArithmeticHead{358}{L}{ANALYSIS.\quad NACHLASS.}
\NumberedItem{5}
\noindent\hspace*{\parindent}Es kommt nur darauf an zu zeigen, dass für die Reihe[\textsuperscript{*)}]
\[
\begin{aligned}
&1-\frac{k(1-\theta)}{n+1}+\frac{k{.}\,k-1}{1{.}\,2}\,\frac{1-\theta{.}\,2-\theta}{n+1{.}\,n+2}\\[.7em]
&\qquad-\frac{k{.}\,k-1{.}\,k-2}{1{.}\,2{.}\,3}\,\frac{1-\theta{.}\,2-\theta{.}\,3-\theta}{n+1{.}\,n+2{.}\,n+3}+\text{u.s.w.},
\end{aligned}
\]
\noindent{\small[S. 292]}\par
\noindent wo $k$ eine ganze positive Zahl ist, eben so wie $n$, und deren Summe
\[
=\frac{n+\theta{.}\,n+\theta+1{.}\,n+\theta+2\,{.}\,{.}\,{.}\,n+\theta+k-1}{n+1{.}\,n+2{.}\,n+3\,{.}\,{.}\,{.}\,n+k}=S
\]
\noindent ist, bei unvollständiger Summirung bis zum $m$ten Gliede aber $=S_m$ gesetzt wird, allemahl $S-S_m$ dasselbe Zeichen hat, wie das erste weggelassene Glied. Für den Fall, wo $1-\theta$ positiv ist, wird dieser Beweis sehr leicht. Es alterniren nehmlich die Zeichen der Glieder der Reihe
\[
1\quad-a\quad+b\quad-c\quad+d\,{.}\,{.}\,{.}\,{.}\,{.}\,g\,{.}\,{.}\,{.}\,{.}\,{.}\,p
\]
\noindent und absolut genommen nehmen sie zu bis zu einem gewissen Gliede z. B. $g$ und von da an nehmen sie ab. Es ist dann klar, dass
\[
\begin{array}{l}
1-S\ \text{positiv und kleiner als }1,\\[.3em]
1-a-S\ \text{negativ und absolut kleiner als }a,\\[.3em]
1-a+b-S\ \text{positiv und kleiner als }b,\\[.3em]
1-a+b-c-S\ \text{negativ u.s.w.}
\end{array}
\]
\noindent sein wird. Die Schlussweise gilt, bis man bei der Summation zu $g$ gelangt ist. Für die übrigen braucht man nur von hinten anzufangen, wo
\begin{center}
\begin{tabular}{l}
das letzte Glied,\\
die algebraische Summe der beiden letzten,\\
die algebraische Summe der drei letzten, u.s.w.
\end{tabular}
\end{center}
\noindent eine Reihe mit alternirenden Zeichen bildet.
\par\smallskip\noindent{\small[S. 293]}\par
\noindent\hspace*{\parindent}Es scheint, dass der Beweis sich viel leichter dadurch führen lässt, dass man nachweiset, aus der Gültigkeit des Satzes bis zu einem gewissen Werth von $k$ folge auch die Richtigkeit für den nächstfolgenden.
\par\medskip\noindent\rule{.25\linewidth}{.4pt}\par\smallskip
{\small\noindent [*) In der Handschrift steht »dass wenn die Reihe«.]\par}
\clearpage
